\begin{sol}[Cilindro rodante en la cavidad de un bloque móvil]

\subsubsection*{Pregunta 1--2: Vínculo cinemático y grados de libertad}

El sistema posee inicialmente tres coordenadas: $x$ (posición del bloque $M$), $\theta$
(posición angular del centro de $m$ respecto a $O$) y $\phi$ (rotación propia de $m$). La
rodadura sin deslizamiento vincula $\theta$ y $\phi$. Cuando el centro de $m$ avanza un
ángulo $\theta$, su centro recorre un arco $(R-r)\theta$ respecto al bloque; igualándolo al
arco de rotación propia $r\phi$:
$$r\phi = (R-r)\theta \;\Longrightarrow\; \dot\phi = \frac{R-r}{r}\dot\theta.$$
La velocidad angular absoluta del cuerpo es
$$\omega = \dot\theta + \dot\phi = \frac{R}{r}\dot\theta.$$
El sistema se reduce así a \textbf{dos grados de libertad}: $x$ y $\theta$.

\subsubsection*{Pregunta 3: Lagrangiano}

Sea $x$ la posición horizontal del centro $O$ de la cavidad. El bloque se mueve solo
horizontalmente ($v_M=\dot x$). El centro de masa de $m$ es
$$\vec{r}_m = \big(x + (R-r)\sin\theta\big)\hat{\imath} - (R-r)\cos\theta\,\hat{\jmath},$$
con
$$v_m^2 = \dot{x}^2 + (R-r)^2\dot\theta^2 + 2\dot{x}(R-r)\dot\theta\cos\theta.$$

Con $I=\gamma m r^2$ ($\gamma=\tfrac12$ cilindro, $\gamma=\tfrac25$ esfera), la energía
cinética es
$$T = \tfrac12 M\dot{x}^2 + \tfrac12 m v_m^2 + \tfrac12 I\omega^2
= \tfrac12(M+m)\dot{x}^2
+ \tfrac12\Big[m(R-r)^2 + I\tfrac{R^2}{r^2}\Big]\dot\theta^2
+ m(R-r)\dot{x}\dot\theta\cos\theta.$$

La energía potencial (nivel de referencia en $O$) es $V=-mg(R-r)\cos\theta$, de modo que
$$\boxed{\;\mathcal{L} = \tfrac12(M+m)\dot{x}^2
+ \tfrac12\Big[m(R-r)^2 + I\tfrac{R^2}{r^2}\Big]\dot\theta^2
+ m(R-r)\dot{x}\dot\theta\cos\theta + mg(R-r)\cos\theta\;}$$

\subsubsection*{Pregunta 4: Ecuaciones de movimiento}

La coordenada $x$ es cíclica ($\partial\mathcal{L}/\partial x=0$), por lo que se conserva
el momento conjugado $p_x=(M+m)\dot{x}+m(R-r)\dot\theta\cos\theta$. Derivando y aplicando
Euler-Lagrange en $\theta$:
$$\boxed{\begin{cases}
(M+m)\ddot{x} + m(R-r)\big(\ddot\theta\cos\theta - \dot\theta^2\sin\theta\big) = 0 \\[4pt]
\Big[m(R-r)^2 + I\left(\dfrac{R}{r}\right)^2\Big]\ddot\theta + m(R-r)\ddot{x}\cos\theta
+ mg(R-r)\sin\theta = 0
\end{cases}}$$
Se sustituye $I=\tfrac12 mr^2$ (cilindro) o $I=\tfrac25 mr^2$ (esfera) según corresponda.

\end{sol}
